% SPDX-FileCopyrightText: 2026 Yiyu Zhou <yiyu@yiyuzhou.io>
%
% SPDX-License-Identifier: GPL-3.0-or-later

\chapter{Making a cue}
\label{ch:cues}

This chapter creates a show of your own next to the storm show, with a cue
that plays one of its sounds.  Besides audio cues, cinq provides video and
picture cues, fade, start, stop, pause, devamp, and GoTo cues, which act on
other cues, group cues, which fire the cues they hold, light and network cues,
and wait and memo cues; later chapters cover each of them.

\section{A new show}

Close the main window to quit cinq, then start it again.  It opens a new show
with an empty cue list, named \texttt{untitled} until you save it.

\section{Adding a cue}

\begin{enumerate}
\item Press \ui{Add cue} in the toolbar.  The drop-down list labeled
  \ui{New cue} sets the kind of cue to add, \ui{Audio} by default.  The new
  cue has no file yet, so its state reads \texttt{NO FILE}, and the sidebar
  shows a warning, as in Figure~\ref{fig:added}.  The inspector shows the
  \ui{Basics} tab for it, whichever tab was open, and the status bar names the
  key that renames the cue.
\item Press \ui{Choose} next to the \ui{File} field of the inspector.  With
  the keys, press \key{Tab} from the cue list, which goes first to the
  fields of the inspector, until the \ui{File} field has the keyboard, and
  press \key{Enter} in the empty field.  Then
  choose \texttt{rain.flac} in the \texttt{sounds} directory of your copy of
  the demo.  The cue takes the file's name, \texttt{rain}, and the warning
  disappears.
\item Press \ui{GO}\@.  The rain plays once, for eight seconds.
\end{enumerate}

\begin{figure}[H]
  \centering
  \includegraphics[width=\textwidth]{screenshots/added.png}
  \caption{New show, with an audio cue that has no file yet}
  \label{fig:added}
\end{figure}

GO skips a cue that cannot play, and clicking its warning selects it.  The
sidebar lists three warnings, or four where the fourth is the last, and counts
the rest.  The
\ui{File} field also accepts a typed path, either absolute or relative to the
show file's directory.  The path applies once you press \key{Enter} or leave
the field, and then the cue takes the file's name, as it does with the
chooser.

\section{Setting the level}

The tutorial says to click a tab of the inspector.  With the keys, \key{gt}
moves to the next tab, and a count before it, such as \key{2gt}, moves to the
tab of that number, counted from one: \key{1gt} shows \ui{Basics}, \key{2gt}
\ui{Time \& Loops}, \key{3gt} \ui{Levels}, and \key{4gt} \ui{Light Patch}\@.
The tab stays where
it is when the selection moves to another cue, so a cue with nothing on that
tab shows a line that says so until you choose another tab.

\begin{enumerate}
\item Click the \ui{Levels} tab of the inspector.
\item Drag the slider to $-20$~dB, or type \texttt{-20} in the field beside
  it and press \key{Enter}\@.
\item Press \ui{Standby selected}, then \ui{GO}\@.  The rain plays more
  quietly.  Move the slider while it plays: it applies as it moves.
\end{enumerate}

A level that you type applies once you press \key{Enter} or leave the field,
so that a cue that plays does not follow each key.

Levels range from $-60$~dB, which is silent, to $+12$~dB\@.  At the default
of 0~dB, the cue plays its file unchanged.

\section{Looping and trimming}

\begin{enumerate}
\item Click the \ui{Time \& Loops} tab, and check \ui{Loop until stopped or
  devamped}.  With the keys, press \key{2gt}, press \key{Tab} until the check box
  has the keyboard, which is the last stop of the tab, and press \key{Enter}
  to check it.  \key{Enter} clears it again.
\item Press \ui{Standby selected}, then \ui{GO}\@.  The rain restarts each
  time it reaches its end.
\item Press the rain's \ui{Stop} button in the sidebar.
\end{enumerate}

\ui{Start} and \ui{End}, on the same tab, limit playback to part of the file,
in seconds from its beginning, which a field takes as minutes and seconds
too, as \texttt{1:30}.  Text that is no time, as \texttt{5s}, shows a red
sentence in the header of the inspector, which says how cinq reads it.  In the
storm show, cue~7 plays the thunder
from 0.2 to 3.5 seconds, and \ui{Duration} shows the resulting length, as in
Figure~\ref{fig:time}.  \ui{Loop start} and \ui{Loop end} set a part of the
file to loop, which Chapter~\ref{ch:sequences} covers.

\begin{figure}[H]
  \centering
  \includegraphics[width=\textwidth]{screenshots/time.png}
  \caption{\ui{Time \& Loops} tab of an audio cue that plays part of its
    file}
  \label{fig:time}
\end{figure}

\section{Editing the list}

The \ui{Basics} tab holds the cue's \ui{Number}, which can be any text, such
as \texttt{1}, \texttt{2.5}, or \texttt{E1}; its \ui{Name}; and its
\ui{Notes}, which the transport shows while the cue is on standby.
\key{Alt+ArrowUp} and \key{Alt+ArrowDown} move the selected cue up and down
the list, and \ui{Delete cue} deletes it.  Cues can also be copied, cut, and
pasted, with the keys that Chapter~\ref{ch:keys} lists.  cinq keeps its own
clipboard, which the desktop does not see.

The status bar at the bottom of the window shows each change along with the
key that undoes it, and after an undo, the key that redoes it.  cinq never
asks for confirmation, since undo is the way back.

Press \key{Control+s}, and save the show as \texttt{my-show.cinq.toml} in
your copy of the demo, where Chapter~\ref{ch:fades} opens it again.

\section*{Exercise}

You can now make a cue that plays a sound once or in a loop, at the level of
your choice.  Add a second audio cue for \texttt{sounds/thunder.wav} that
plays only its first two seconds.
